\documentclass[10pt,twocolumn]{article}
\usepackage[utf8]{inputenc}\usepackage[T1]{fontenc}\usepackage[ngerman]{babel}\usepackage{amsmath,amssymb}\usepackage[margin=1.8cm]{geometry}
\usepackage{times}\usepackage{hyperref}
\title{\textbf{Spektren, Rotation und Breite: ein verifikator-gesteuertes Agentenlabor zu Optimierern und Skalierungsgesetzen}}
\author{Verifier-Gated Discovery Lab \and Team Ninja Turtles\\ \small Hack-Nation 2026, Challenge 3}
\date{Preprint, \today}
\begin{document}\maketitle
\section{1 Einleitung}

\textbf{Frage.} Warum funktionieren moderne Optimierer und Lernraten-Schedules, wann überträgt sich die Lernrate über die Modellbreite, und woher stammen die Exponenten von Skalierungsgesetzen? Diese Arbeit berichtet, was ein verifikator-gesteuertes Agentenlabor darauf numerisch belegen konnte und was nicht.

\textbf{Kontext.} Die Literaturlage im Labor ist folgende:
\begin{itemize}
\item Das spektrale Schrittweitenprofil (volatiler Kopf an der Stabilitätsgrenze, toleranter Bulk) wird als Erklärung dafür angeführt, dass Muon Adam und Adam SGD übertrifft {\scriptsize[C-lit3]} (arXiv:2608.25990v1).
\item Im linearen Assoziativspeicher-Modell mit Potenzgesetz-Spektrum wird Muons Skalierungsgesetz hergeleitet und seine bessere Skalierungseffizienz gegenüber GD gezeigt {\scriptsize[C-lit8]} (arXiv:2602.05725v3).
\item Unter μP bleiben viele optimale Hyperparameter über Modellgrößen stabil {\scriptsize[C-lit29]} (arXiv:2203.03466v2).
\end{itemize}

\textbf{Beitrag.} Das Labor prüft solche Aussagen in numerischen Modellwelten. Eine Behauptung gilt nur, wenn ein unabhängiger Code-Prüfer sie mit eigenen Seeds, eigenem Lernraten-Tuning und eigener Auflösung nachrechnet. Red-Team-Gegenprüfungen versuchen jede Aussage zu widerlegen.

\textbf{Zusammenfassung der Resultate.}
\begin{itemize}
\item Gradientenfluss: Die Verlustexponenten stimmen an vier geprüften (a, b)-Paaren mit den behaupteten Werten überein (Befund I) {\scriptsize[C-optscale-R17]}.
\item Rechenoptimum bei C = P·t: Es wurden sieben (a, b)-Punkte bestätigt, und die Alternative "reiner Gradientenfluss-Exponent" wurde nicht gezeigt (Befund II) {\scriptsize[C-optscale-R6]} {\scriptsize[C-optscale-R8]} {\scriptsize[C-optscale-R9]}.
\item Random-Feature- und Ridge-Exponenten sind nur an wenigen Punkten bestätigt. In Regimen mit kleinem a und großem b liegen die Fits unter dem Abschneide-Exponenten, und die asymptotischen Exponenten sind ungeklärt (Befunde III und IV) {\scriptsize[C-optscale-R1-RT1]} {\scriptsize[C-optscale-R1-RT2]}.
\item Muon gegen Adam hängt von Rotation und Spektrum ab. Rotation allein entscheidet das Vorzeichen nicht (Beobachtungen V und VI) {\scriptsize[C-optscale-R2]} {\scriptsize[C-optscale-R4]} {\scriptsize[C-optscale-R5]}.
\item Unter heavy-tailed Rauschen schlagen Shampoo und Muon SGD (Beobachtung VII) {\scriptsize[C-optscale-R13]}.
\item Lernraten-Schedules senken den Endverlust in den geprüften Fällen (Beobachtung VIII) {\scriptsize[C-optscale-R15]}.
\item Edge of Stability tritt bei einer geprüften Lernrate in allen fünf Seeds auf (Befund IX) {\scriptsize[C-optscale-R16]}.
\item Für GD und Heavy-Ball auf einer rationalen Quadratik liegen exakte Stabilitätszertifikate vor (Proposition 1) {\scriptsize[C-optscale-R18]}.
\item μP überträgt die optimale Lernrate über Breiten, SP nicht (Befund X) {\scriptsize[C-optscale-R10]}.
\end{itemize}

\section{2 Modell}

Die Annahmen, die die Resultate tragen, stehen jeweils in den Prüfaufträgen der Claims.

\begin{itemize}
\item \textbf{Lineares Modell mit Potenzgesetz-Spektrum.} Die Eigenwerte sind lambda_k = k^-a und die Zielanteile k^-b {\scriptsize[C-optscale-R1]} {\scriptsize[C-optscale-R17]}. Betrachtet werden der Gradientenfluss-Verlust über die Zeit t, der Verlust des linearen Random-Feature-Modells über P {\scriptsize[C-optscale-R1]}, der Exzess-Risiko-Verlauf der Ridge-Regression über N {\scriptsize[C-optscale-R12]} und das rechenoptimale Budget C = P·t {\scriptsize[C-optscale-R6]}. Ein Exponent e bedeutet Verlust ~ x^-e.
\item \textbf{Optimierer auf Quadratiken.} Es geht um Kronecker-Krümmung mit Spektren a_L und a_R, wahlweise rotiert. Verglichen werden adam, muon, shampoo, sgd, signum und weitere mit Schedules (const, cosine, wsd) bei selbst getunter Lernrate {\scriptsize[C-optscale-R2]} {\scriptsize[C-optscale-R13]} {\scriptsize[C-optscale-R15]}. Die Metrik ist der relative Endverlust.
\item \textbf{Edge of Stability.} Voll-Batch-GD auf einem tanh-Netz, Kenngröße lr·lambda_max/2 {\scriptsize[C-optscale-R16]}.
\item \textbf{μP gegen SP.} Es wird geprüft, ob sich die optimale Lernrate mit der Breite verschiebt {\scriptsize[C-optscale-R10]}.
\item \textbf{Exakte Anker.} GD und Heavy-Ball auf Quadratiken mit rationaler Matrix H, mit der Stabilitätsbedingung |beta| < 1, H > 0 und H < 2(1+beta)/lr {\scriptsize[C-optscale-R18]}.
\end{itemize}

\section{3 Methode: das agentische Labor}

Das Labor lief 18 Runden, 18 geprüfte Aussagen, 0 negative Ergebnisse, Kosten 8.86 USD {\scriptsize[C-methode]}. Jede Runde wurde vor dem Experiment präregistriert (prereg.md) {\scriptsize[C-methode]}.

Die Claims dokumentieren drei Mechanismen:
1. \textbf{Code-Prüfer.} Er rechnet mit eigenen Seeds nach, tunet die Lernrate selbst und fittet Exponenten in zwei Fenstern mit fester Toleranz. Für exakte Anker verwendet er rationale Arithmetik {\scriptsize[C-optscale-R17]} {\scriptsize[C-optscale-R18]}.
2. \textbf{Red-Team.} Es formuliert Gegenprüfungen zu jeder Aussage. Eine Aussage gilt als \emph{angefochten}, sobald eine Gegenprüfung besteht, auch wenn diese mit der Aussage vereinbar ist. Angefochten sind optscale-R7, optscale-R15 und optscale-R18. Im Einzelfall ist zu lesen, ob die bestandene Gegenprüfung der Aussage widerspricht {\scriptsize[C-redteam-methode]}.
3. \textbf{Präregistrierte Mehrfachtest-Korrektur.} Benjamini-Hochberg mit q = 0.1 über alle 20 Optimierer-Vergleiche des Prüfers {\scriptsize[C-bh]}. Prüfungen mit Lernraten-Optimum am Gitterrand gelten als nicht entscheidbar und zählen nicht als Test {\scriptsize[C-bh]}.

Die Rollen Scout, Integrator und Forscher sind in der Claim-Liste nicht beschrieben und werden hier nicht weiter ausgeführt.

\textbf{Lesregel.} Eine nicht bestandene Gegenprüfung zeigt nur, dass ihre Behauptung nicht gezeigt ist, nicht das Gegenteil. Die allgemeine Formel hinter mehreren geprüften Einzelpunkten ist eine Interpretation, kein Theorem.

\section{4 Resultate}

\subsubsection*{4.1 Skalierung im linearen Modell}

\textbf{Numerischer Befund I (observed).} Der Gradientenfluss-Verlust über t fällt an vier Punkten mit den Prüfer-Exponenten in beiden Fenstern innerhalb der Toleranz 0.03 wie behauptet {\scriptsize[C-optscale-R17]}:

| (a, b) | Exponent | Beleg |
|---|---|---|
| (1.0, 2.0) | 1.0 | {\scriptsize[C-optscale-R17]} |
| (1.0, 3.0) | 2.0 | {\scriptsize[C-optscale-R17]} |
| (2.0, 4.0) | 1.5 | {\scriptsize[C-optscale-R17]} |
| (0.5, 1.5) | 1.0 | {\scriptsize[C-optscale-R17]} |

Zusätzlich lieferte der Prüfer bei (a=2.0, b=3.0) den Exponenten 1.0000 {\scriptsize[C-optscale-R17-RT1]} und bei (a=0.5, b=3.0) den Exponenten 4.0000 {\scriptsize[C-optscale-R17-RT2]}. Beide Werte liegen in beiden Fenstern vor. Die Alternative e = (b−1)·a wurde nicht gezeigt {\scriptsize[C-optscale-R17-RT1]} {\scriptsize[C-optscale-R17-RT2]}.

\textbf{Numerischer Befund II (observed).} Der rechenoptimale Verlust bei Budget C = P·t fällt an folgenden Punkten wie behauptet, mit Toleranz 0.03 {\scriptsize[C-optscale-R6]} {\scriptsize[C-optscale-R8]} {\scriptsize[C-optscale-R9]}:

| (a, b) | Exponent | Beleg |
|---|---|---|
| (1.0, 2.0) | 0.5 | {\scriptsize[C-optscale-R6]} |
| (2.0, 3.0) | 0.6667 | {\scriptsize[C-optscale-R6]} |
| (1.0, 4.0) | 1.5 | {\scriptsize[C-optscale-R6]} |
| (1.5, 1.5) | 0.2 | {\scriptsize[C-optscale-R6]} |
| (0.5, 2.0) | 0.6667 | {\scriptsize[C-optscale-R8]} |
| (1.0, 3.0) | 1.0 | {\scriptsize[C-optscale-R8]} |
| (2.0, 2.0) | 0.3333333333 | {\scriptsize[C-optscale-R9]} |

Die Gegenhypothese des reinen Gradientenfluss-Exponenten (b−1)/a wurde nicht gezeigt. Bei (a=2.0, b=3.0) lieferte der Prüfer 0.6658 und 0.6663 statt 1.0 {\scriptsize[C-optscale-R6-RT1]}. Bei (a=3.0, b=4.0) lieferte er 0.7474 und 0.7484 statt 1.0 {\scriptsize[C-optscale-R6-RT2]}.

Die Formel C^-(b−1)/(a+1) fasst diese Punkte zusammen. Sie ist eine Interpretation {\scriptsize[C-optscale-R6-I]} und kein Resultat.

Die Claims R8 und R9 fragen nach dem Budget C = P·N mit Random-Feature und Ridge. Die vom Prüfer geprüfte Kurve ist jedoch die Gradientenfluss-Kurve "compute", nicht die RF-Ridge-Kurve {\scriptsize[C-optscale-R8-RT2]}. Auch R7 stellt die Frage nach einem RF-basierten Optimum. Bestätigt ist dort nur der RF-Exponent über P bei (1.0, 2.0) (Befund III), während die bestandenen Gegenprüfungen wiederum die Gradientenfluss-Kurve betreffen {\scriptsize[C-optscale-R7]} {\scriptsize[C-optscale-R7-RT1]} {\scriptsize[C-optscale-R7-RT2]}. Der Exponent für C = P·N bleibt daher ungeprüft (siehe die Abschnitte zu negativen Ergebnissen und zu Grenzen).

\textbf{Numerischer Befund III (observed).} Der Verlust des linearen Random-Feature-Modells über P fällt an zwei Punkten wie x^-1.0 (Toleranz 0.1) {\scriptsize[C-optscale-R1]} {\scriptsize[C-optscale-R11]}:
\begin{itemize}
\item bei (a=1.0, b=2.0) mit Exponenten 0.9909 und 1.0540 {\scriptsize[C-optscale-R1]};
\item bei (a=2.0, b=2.0) mit Exponenten 1.0782 und 1.0806 {\scriptsize[C-optscale-R11]}.
\end{itemize}

\textbf{Numerischer Befund IV (observed).} Der Exzess-Risiko-Verlust der Ridge-Regression über N (a=1.0, b=3.0, ridge = 0.001, noise = 0) fällt wie x^-2.0, mit Exponenten 1.9893 und 1.9258 (Toleranz 0.1) {\scriptsize[C-optscale-R12]}.

\subsubsection*{4.2 Optimierer auf Quadratiken}

\textbf{Beobachtung V (statistical).} Ohne Rotation bei (a_L, a_R) = (3.0, 3.0) und 500 Schritten erreicht Adam (const) einen kleineren relativen Endverlust als Muon (const) {\scriptsize[C-optscale-R2]} {\scriptsize[C-optscale-R3]}. Die Werte sind 4.932e-23 gegen 2.825e-04, mit p = 0.0008 auf 10 frischen Seeds und eigenem Tuning (beste Lernraten 2^-2 und 2^-5) {\scriptsize[C-optscale-R2]} {\scriptsize[C-optscale-R3]}.

\textbf{Beobachtung VI (statistical).} Mit Rotation hängt das Vorzeichen vom Spektrum ab {\scriptsize[C-optscale-R4]} {\scriptsize[C-optscale-R5]} {\scriptsize[C-optscale-R14]}:

| (a_L, a_R) | Befund | Verlust Muon / Adam | Verhältnis | Beleg |
|---|---|---|---|---|
| (3.0, 3.0) | Muon besser | 4.481e-04 / 1.810e-03 | 0.248 (p = 0.0008) | {\scriptsize[C-optscale-R4]} |
| (2.5, 2.5) | Muon besser | 1.814e-04 / 1.683e-03 | 0.108 (p = 0.0008) | {\scriptsize[C-optscale-R14]} |
| (1.0, 1.0) | Adam besser | 8.863e-05 / 3.170e-11 | p = 0.0041 | {\scriptsize[C-optscale-R4]} {\scriptsize[C-optscale-R5]} |

Rotation allein entscheidet das Vorzeichen daher nicht. Bei (3.0, 3.0) kehrt es sich zwischen unrotiert und rotiert um {\scriptsize[C-optscale-R2]} {\scriptsize[C-optscale-R4]}. Bei (1.0, 1.0) rotiert gewinnt Adam {\scriptsize[C-optscale-R5]}. Die geprüften Punkte begrenzen den Wechsel auf das Intervall zwischen a = 1.0 und a = 2.5 {\scriptsize[C-optscale-R5]} {\scriptsize[C-optscale-R14]}. Seine genaue Lage ist nicht bestätigt.

Der Befund bei (1.0, 1.0) steht in den Claims R4 und R5. Der BH-Eintrag zu R5 nennt zudem ein Verhältnis 0.107 für Muon gegen Adam, ohne den zugehörigen Punkt auszuweisen {\scriptsize[C-optscale-R5]} {\scriptsize[C-bh]}.

\textbf{Beobachtung VII (statistical).} Unter Student-Rauschen (df = 2, sigma = 0.1) auf der rotierten Quadratik mit (a_L, a_R) = (2.0, 2.0) erreichen Shampoo und Muon kleinere Endverluste als SGD {\scriptsize[C-optscale-R13]}:

| Vergleich | Verlust | Verhältnis | Beleg |
|---|---|---|---|
| Shampoo gegen SGD | 5.032e-02 / 9.663e-02 | 0.521 | {\scriptsize[C-optscale-R13]} |
| Muon gegen SGD | 5.483e-02 / 9.663e-02 | 0.567 | {\scriptsize[C-optscale-R13]} |

Beide haben p = 0.0008 {\scriptsize[C-optscale-R13]}. Die Claims bestätigen keinen Vergleich von Clipping gegen SGD und keinen Vergleich von Adam oder Signum gegen SGD zugunsten der anderen Seite (siehe den Abschnitt zu negativen Ergebnissen).

\subsubsection*{4.3 Schedules}

\textbf{Beobachtung VIII (statistical, angefochten).} Auf der rotierten Quadratik mit (2.5, 2.5) und Gauss-Rauschen (sigma = 0.1) erreicht Adam mit Cosine-Schedule einen kleineren Endverlust als mit const {\scriptsize[C-optscale-R15]}. Die Werte sind 1.142e-02 gegen 3.190e-02, Verhältnis 0.358, p = 0.0008 {\scriptsize[C-optscale-R15]}.

Der Befund ist nicht auf Cosine mit Adam beschränkt. SGD mit WSD schlägt SGD mit const unter Student-Rauschen (df = 3): 1.422e-02 gegen 2.563e-02, Verhältnis 0.555 {\scriptsize[C-optscale-R15-RT2]} {\scriptsize[C-bh]}. Die Anfechtung beruht auf dieser bestandenen Gegenprüfung, die den Befund erweitert, ihm aber nicht widerspricht {\scriptsize[C-redteam-methode]}.

\subsubsection*{4.4 Edge of Stability}

\textbf{Numerischer Befund IX (observed).} Voll-Batch-GD auf dem tanh-Netz (Breite 16, lr = 1.0, 2000 Schritte) hat am Ende lr·lambda_max/2 in [0.85, 1.25], und zwar in 5 von 5 frischen Seeds {\scriptsize[C-optscale-R16]}. Die Werte bei 80 \% und 100 \% der Schritte liegen zwischen 0.958 und 1.152, gemessen am vollen Hesse-Eigenwert {\scriptsize[C-optscale-R16]}.

\textbf{Proposition 1 (computed_rigorous; optscale-R18 ist angefochten).} Exakte Zertifikate in rationaler Arithmetik auf H = diag(1, 2) {\scriptsize[C-optscale-R18]} {\scriptsize[C-optscale-R18-RT1]} {\scriptsize[C-optscale-R18-RT2]}:
\begin{itemize}
\item GD mit lr = 99/100 ist stabil, mit lr = 101/100 nicht {\scriptsize[C-optscale-R18]}.
\item GD mit lr = 1 ist nicht stabil. Das ist der Randfall 2/lr = 2 mit Spektralradius 1 {\scriptsize[C-optscale-R18-RT1]}.
\item Heavy-Ball mit lr = 3/2 und beta = 9/10 ist stabil (Grenze 38/15) {\scriptsize[C-optscale-R18]}.
\item Heavy-Ball mit lr = 19/10 und beta = 9/10 ist nicht stabil {\scriptsize[C-optscale-R18-RT2]}.
\end{itemize}

Die Anfechtung beruht auf den bestandenen Gegenprüfungen zu lr = 1 und lr = 19/10 {\scriptsize[C-optscale-R18-RT1]} {\scriptsize[C-optscale-R18-RT2]}. Diese sind mit der Aussage vereinbar, da sie genau die Randfälle der Bedingung H < 2(1+beta)/lr belegen {\scriptsize[C-redteam-methode]}.

\subsubsection*{4.5 μP gegen SP}

\textbf{Numerischer Befund X (observed).} Optimale log2-Lernraten über die Breiten 32, 64, 128, 256 (Gitter mit Faktor 2, 300 Schritte, 2 Seeds) {\scriptsize[C-optscale-R10]}:

| Parametrisierung, Optimierer | Optima (kleinste bis größte Breite) | Verschiebung | Spannweite | Kriterium | Beleg |
|---|---|---|---|---|---|
| μP, Adam | -7/-6/-7/-7 | 0 | 1 | Transfer | {\scriptsize[C-optscale-R10]} |
| SP, Adam | -6/-6/-7/-8 | -2 | 2 | kein Transfer | {\scriptsize[C-optscale-R10]} |
| μP, SGD | -1/-1/0/-1 | 0 | 1 | Transfer | {\scriptsize[C-optscale-R10]} |
| SP, SGD | 0/-1/-1/-2 | -2 | 2 | kein Transfer | {\scriptsize[C-optscale-R10]} |

Dieser Befund ist verträglich mit dem berichteten μP-Transfer in der Literatur {\scriptsize[C-lit29]} (arXiv:2203.03466v2).

\subsubsection*{4.6 Mehrfachtest-Korrektur}

\textbf{Beobachtung XI (statistical).} Von 20 eindeutigen Optimierer-Vergleichen des Prüfers bleiben nach Benjamini-Hochberg (q = 0.1) 11 signifikant mit Verhältnis unter 0.9 {\scriptsize[C-bh]}. Das sind die Hauptprüfungen von R2, R3, R4, R5, R13, R14 und R15 sowie die Gegenprüfung R15-RT2 {\scriptsize[C-bh]}. Die nicht signifikanten Tests sind Gegenprüfungen mit p_adj = 1.0, außer R4-RT2 mit p = 0.9967 und R13-RT2 mit p = 0.994 {\scriptsize[C-bh]}.

\section{5 Negative Ergebnisse und Red-Team-Befunde}

Nicht bestandene Gegenprüfungen belegen nur, dass ihre Behauptung nicht gezeigt ist.

\textbf{Skalierung (Gegenprüfungen nicht bestanden).}
\begin{itemize}
\item \textbf{Random-Feature-Exponent über P bei kleinem a und großem b.} Die Behauptungen eines RF-Exponenten von 1.0 bei (a=0.5, b=3.5) und von 0.6 bei (a=0.3, b=4.0) sind nicht gezeigt {\scriptsize[C-optscale-R1-RT1]} {\scriptsize[C-optscale-R1-RT2]}. Die Prüfer-Fits sind 0.6189 und 0.7694 bzw. 0.3146 und 0.4394 {\scriptsize[C-optscale-R1-RT1]} {\scriptsize[C-optscale-R1-RT2]}. Auch die Behauptung 1.0 bei (a=0.5, b=3.0) ist nicht gezeigt (0.6077 und 0.7600) {\scriptsize[C-optscale-R11-RT1]}. Alle Fits liegen unter den im Red-Team genannten Abschneide-Exponenten 2,5 bzw. 3 und steigen mit dem Fenster {\scriptsize[C-optscale-R1-RT1]} {\scriptsize[C-optscale-R1-RT2]}. Die Sättigungsformel P^-min(b-1, 2a) ist dadurch nicht bestätigt, und der asymptotische Exponent bleibt offen.
\item \textbf{Ridge-Exponent über N in Sättigungsregimen.} Nicht gezeigt sind 1.0 bei (a=0.5, b=3.0, 1.4644 und 1.8149) {\scriptsize[C-optscale-R12-RT1]}, 0.8 bei (a=0.4, b=4.0, 1.2112 und 1.5759) {\scriptsize[C-optscale-R12-RT2]} und 2.0 bei (a=1.0, b=4.0, 2.5008 und 2.7484) {\scriptsize[C-optscale-R9-RT2]}. Die rauschbehaftete Behauptung 0.5 bei (a=2.0, b=2.0, noise = 1.0) ist ebenfalls nicht gezeigt, die Fits sind -0.1666 und -0.2154 {\scriptsize[C-optscale-R11-RT2]}.
\item \textbf{RF-Exponent bei C = P·N.} Nicht gezeigt sind 0.5 bei (a=2.0, b=2.0, Fits 1.0782 und 1.0806) und 1.0 bei (a=2.0, b=3.0, Fits 1.8814 und 1.8742) {\scriptsize[C-optscale-R8-RT1]} {\scriptsize[C-optscale-R8-RT2]}. Auch der behauptete compute-Exponent 1.2 bei (a=1.0, b=4.0) ist nicht gezeigt; der Prüfer maß dort 1.4990 und 1.4998 {\scriptsize[C-optscale-R9-RT1]}.
\end{itemize}

\textbf{Optimierer (Gegenprüfungen nicht bestanden).}
\begin{itemize}
\item Muon gleich gut oder besser als Adam bei a = 0.1 ohne Rotation ist nicht gezeigt {\scriptsize[C-optscale-R3-RT2]}. Gemessen wurden 1.162e-05 gegen 3.883e-23 {\scriptsize[C-optscale-R3-RT2]}. Auch rotiert bei a = 0.1 ist Muon nicht besser (1.131e-05 gegen 3.873e-23) {\scriptsize[C-optscale-R5-RT2]}.
\item Adam besser als Muon im rotierten Fall mit Ungleichgewicht a_L ≠ a_R ist nicht gezeigt (2.605e-04 gegen 3.093e-05) {\scriptsize[C-optscale-R2-RT2]} {\scriptsize[C-optscale-R3-RT1]}. Dieselben Zahlen stehen in beiden Claims.
\item Muon besser als Adam bei (1.0, 1.0) rotiert ist nicht gezeigt (8.863e-05 gegen 3.170e-11) {\scriptsize[C-optscale-R4-RT2]} {\scriptsize[C-optscale-R5-RT1]}.
\item Adam besser als Muon bei a = 2.5 und a = 3.0 rotiert ist nicht gezeigt {\scriptsize[C-optscale-R14-RT1]} {\scriptsize[C-optscale-R14-RT2]}.
\item SGD besser als Adam oder Signum unter Student-Rauschen ist nicht gezeigt (9.663e-02 gegen 8.435e-02 bzw. 8.455e-02) {\scriptsize[C-optscale-R13-RT1]} {\scriptsize[C-optscale-R13-RT2]}.
\item Cosine besser als const bei SGD im MLP ist nicht gezeigt (4.859e-03 gegen 3.732e-03) {\scriptsize[C-optscale-R15-RT1]}.
\end{itemize}

\textbf{Nicht entscheidbar.} Bei der Gegenprüfung zu Shampoo ohne Rotation und bei Shampoo gegen Adam bei (3, 3) rotiert lag das Lernraten-Optimum am Gitterrand 2^-16..2^4 {\scriptsize[C-optscale-R2-RT1]} {\scriptsize[C-optscale-R4-RT1]}. Sie sind weder bestanden noch widerlegt und zählen nicht als Test {\scriptsize[C-bh]}.

\textbf{EOS (Gegenprüfungen nicht bestanden).} Der EOS-Wert am Ende ist nicht unter Wert 1 gedrückt (Werte 0.958 bis 1.152) {\scriptsize[C-optscale-R16-RT1]}. Dass EOS bei kleinerer Lernrate in 4000 Schritten erreicht wird, ist nicht gezeigt {\scriptsize[C-optscale-R16-RT2]}. Dort liegen die Werte zwischen 0.595 und 0.896 {\scriptsize[C-optscale-R16-RT2]}.

\textbf{μP und SGD (Gegenprüfungen nicht bestanden).} Der Befund "SGD unter SP verschiebt sich nicht" und der Befund "SGD unter μP verschiebt sich um mehr als eine Gitterstufe" sind beide nicht gezeigt {\scriptsize[C-optscale-R10-RT1]} {\scriptsize[C-optscale-R10-RT2]}.

\section{6 Grenzen und offene Fragen}

\begin{itemize}
\item \textbf{Formeln sind Interpretationen.} Die Formeln e = (b−1)/a {\scriptsize[C-optscale-R17-I]} und C^-(b-1)/(a+1) {\scriptsize[C-optscale-R6-I]} sind ungeprüfte Interpretationen mehrerer geprüfter Einzelpunkte, keine Theoreme. Ebenso ungeprüft sind die Aufteilung P_opt ~ C^(1/(1+a)), t_opt ~ C^(a/(1+a)) {\scriptsize[C-optscale-R6-I]} und die vermutete Sättigung der Random-Feature- und Ridge-Exponenten bei 2a {\scriptsize[C-optscale-R12-I]}.
\item \textbf{Rechenoptimum bei C = P·N.} Der Exponent für Random-Feature plus optimierte Ridge ist nicht bestätigt. Das gilt auch für die Zerlegung in e_P und e_N {\scriptsize[C-optscale-R9]} {\scriptsize[C-optscale-R11-I]}. Die Fits in Regimen mit kleinem a und großem b sind fensterabhängig (siehe den Abschnitt zu negativen Ergebnissen) {\scriptsize[C-optscale-R1-RT1]} {\scriptsize[C-optscale-R1-RT2]}.
\item \textbf{Lage und Art des Muon-Adam-Wechsels.} Die genaue Lage zwischen a = 1.0 und a = 2.5 und die Frage "Phasenübergang oder glatter Übergang" sind offen {\scriptsize[C-optscale-R5]} {\scriptsize[C-optscale-R14]} {\scriptsize[C-optscale-R4-I]}. Die Rotation ist nur binär steuerbar, graduelle Rotation wurde nicht getestet {\scriptsize[C-optscale-R4-I]}. Die Hypothese "Rotation sei der alleinige Treiber" {\scriptsize[C-optscale-R3-I]} ist durch die geprüften Punkte mit Rotation und Adam-Vorteil bei (1.0, 1.0) nicht gedeckt {\scriptsize[C-optscale-R5]}. Die Hypothese eines Wechsels zwischen 1.0 und 1.5 {\scriptsize[C-optscale-R5-I]} ist ungeprüft.
\item \textbf{Shampoo.} Der Vergleich mit Adam ohne und mit Rotation ist wegen Randoptima nicht entscheidbar {\scriptsize[C-optscale-R2-RT1]} {\scriptsize[C-optscale-R4-RT1]}. Zu Clipping, Signum und Adam unter heavy-tailed Rauschen liegt kein bestätigter Vorteil gegenüber SGD vor {\scriptsize[C-optscale-R13]}. Die Literatur zu heavy-tailed Konvergenzraten {\scriptsize[C-lit18]} (arXiv:2602.07425v2) ist nicht mit eigenen Messungen verknüpft.
\item \textbf{Schedules.} Bestätigt sind nur Adam/cosine bei Gauss-Rauschen und SGD/wsd bei Student-Rauschen mit df = 3 {\scriptsize[C-optscale-R15]} {\scriptsize[C-optscale-R15-RT2]}. Dass Schedules im MLP helfen, ist nicht gezeigt {\scriptsize[C-optscale-R15-RT1]}.
\item \textbf{EOS.} Die Schwelle der Lernrate, ab der EOS auftritt, ist nicht bestimmt. Gezeigt ist nur ein Punkt bei Breite 16 {\scriptsize[C-optscale-R16]}.
\item \textbf{μP gegen SP.} Der Transfer ist nur an vier Breiten mit grobem Faktor-2-Gitter, wenigen Schritten und 2 Seeds geprüft. Die Verschiebung pro Breitenverdopplung ist nicht quantifiziert {\scriptsize[C-optscale-R10]} {\scriptsize[C-optscale-R10-I]}.
\item \textbf{Exakte Anker.} Die Zertifikate betreffen nur diagonale rationale Matrizen und einzelne Lernraten {\scriptsize[C-optscale-R18]}.
\item \textbf{Statistik.} Die BH-Korrektur deckt nur die 20 Optimierer-Vergleiche ab, nicht die Exponenten-Fits {\scriptsize[C-bh]}.
\end{itemize}

\section{Literatur}

Nur Quellen aus Claims:
\begin{itemize}
\item arXiv:2203.03466v2 {\scriptsize[C-lit29]} {\scriptsize[C-lit30]}
\item arXiv:2602.05725v3 {\scriptsize[C-lit8]} {\scriptsize[C-lit9]}
\item arXiv:2602.07425v2 {\scriptsize[C-lit18]}
\item arXiv:2608.25990v1 {\scriptsize[C-lit3]} {\scriptsize[C-lit4]}
\end{itemize}

---
Prüfprotokoll: 72 Claims zitiert, Korrekturrunden [{"runde": 0, "verstoesse": 54}, {"runde": 1, "verstoesse": 8}, {"final_verstoesse_entfernt": 0}], 0 unbelegte Sätze entfernt, verbleibende Verstöße: 0.
\end{document}
